\noindent
\noindent
La presente sección del anexo se fundamenta en el trabajo de grado desarrollado por Michael Stiven Honores Quishpillo y Juan David Naranjo Sánchez, titulado ``Propuesta de Transformación Digital para la Gestión de la Información en el Proceso de Producción Establecido en la Organización Prefabricados JAMAR Validada a través de la Construcción de un Prototipo Funcional''.

\section{Elementos de negocio}

\begin{longtable}{|c|p{8cm}|}
	\caption{Elementos de negocio en ArchiMate} \label{tab:elementos-negocio-archimate}                                                  \\
	\hline
	\textbf{Icono}                                             & \textbf{Descripción}                                                    \\
	\hline
	\endfirsthead

	\caption[]{Elementos de negocio en ArchiMate (continuación)}                                                                         \\
	\hline
	\textbf{Icono}                                             & \textbf{Descripción}                                                    \\
	\hline
	\endhead

	\hline
	\endfoot

	\endlastfoot
	\includegraphics{anexos/ARCHI/business/actor.png}          &
	\textbf{Business Actor:} Corresponde a una entidad capaz de ejecutar comportamientos dentro del contexto organizacional, como puede ser una persona o una organización. \\
	\hline
	\includegraphics{anexos/ARCHI/business/rol.png}            &
	\textbf{Business Role:} Representa el conjunto de responsabilidades y comportamientos que pueden ser asumidos por un actor dentro del negocio. \\
	\hline
	\includegraphics{anexos/ARCHI/business/colaboration.png}   &
	\textbf{Business Collaboration:} Se refiere a la agrupación de dos o más roles de negocio que cooperan para alcanzar un objetivo común. \\
	\hline
	\includegraphics{anexos/ARCHI/business/interface.png}      &
	\textbf{Business Interface:} Constituye un punto de acceso mediante el cual un servicio de negocio es ofrecido a los actores. \\
	\hline
	\includegraphics{anexos/ARCHI/business/process.png}        &
	\textbf{Business Process:} Conjunto estructurado de actividades orientadas a la obtención de un resultado específico para un cliente interno o externo. \\
	\hline
	\includegraphics{anexos/ARCHI/business/function.png}       &
	\textbf{Business Function:} Agrupación de comportamientos que comparten una misma base de recursos dentro del negocio. \\
	\hline
	\includegraphics{anexos/ARCHI/business/interaction.png}    &
	\textbf{Business Interaction:} Describe un comportamiento que surge de la interacción entre dos o más roles de negocio. \\
	\hline
	\includegraphics{anexos/ARCHI/business/service.png}        &
	\textbf{Business Service:} Representa un servicio que satisface una necesidad de un cliente, ya sea interno o externo a la organización. \\
	\hline
	\includegraphics{anexos/ARCHI/business/event.png}          &
	\textbf{Business Event:} Hace referencia a un suceso que impacta la continuidad o el flujo de un proceso de negocio. \\
	\hline
	\includegraphics{anexos/ARCHI/business/object.png}         &
	\textbf{Business Object:} Concepto utilizado dentro de un dominio específico del negocio. \\
	\hline
	\includegraphics{anexos/ARCHI/business/contract.png}       &
	\textbf{Contract:} Acuerdo, ya sea formal o informal, que define los derechos y obligaciones asociados a un producto. \\
	\hline
	\includegraphics{anexos/ARCHI/business/representation.png} &
	\textbf{Representation:} Forma en la que la información contenida en un objeto de negocio se presenta o es percibida. \\
	\hline
	\includegraphics{anexos/ARCHI/business/product.png}        &
	\textbf{Product:} Conjunto coherente de servicios y/o objetos pasivos, acompañado de acuerdos o contratos que lo respaldan. \\
	\hline
	\includegraphics{anexos/ARCHI/business/grouping.png}       &
	\textbf{Grouping:} Mecanismo utilizado para organizar y agrupar elementos relacionados dentro de un modelo de arquitectura. \\                                  \\
\end{longtable}

\section{Elementos de aplicación}

\begin{longtable}{|c|p{8cm}|}
	\caption{Elementos de aplicación en ArchiMate} \label{tab:elementos-aplicacion-archimate}                                                            \\
	\hline
	\textbf{Icono}                                               & \textbf{Descripción}                                                                  \\
	\hline
	\endfirsthead

	\caption[]{Elementos de aplicación en ArchiMate (continuación)}                                                                                      \\
	\hline
	\textbf{Icono}                                               & \textbf{Descripción}                                                                  \\
	\hline
	\endhead

	\hline
	\endfoot

	\endlastfoot
	\includegraphics{anexos/ARCHI/application/component.png}     &
	\textbf{Application Component:} Unidad modular que proporciona un conjunto de servicios relacionados dentro del ámbito de las aplicaciones. \\
	\hline
	\includegraphics{anexos/ARCHI/application/collaboration.png} &
	\textbf{Application Collaboration:} Interacción establecida entre dos o más componentes de aplicación que cooperan entre sí. \\
	\hline
	\includegraphics{anexos/ARCHI/application/interface.png}     &
	\textbf{Application Interface:} Punto de acceso a las funcionalidades que ofrece una aplicación. \\
	\hline
	\includegraphics{anexos/ARCHI/application/process.png}       &
	\textbf{Application Process:} Conjunto de comportamientos ejecutados por uno o varios componentes de aplicación con el fin de alcanzar un objetivo determinado. \\
	\hline
	\includegraphics{anexos/ARCHI/application/function.png}      &
	\textbf{Application Function:} Agrupación de comportamientos relacionados en el contexto de las aplicaciones. \\
	\hline
	\includegraphics{anexos/ARCHI/application/interaction.png}   &
	\textbf{Application Interaction:} Describe el comportamiento resultante de la interacción entre dos o más componentes de aplicación. \\
	\hline
	\includegraphics{anexos/ARCHI/application/service.png}       &
	\textbf{Application Service:} Servicio que una aplicación expone hacia su entorno. \\
	\hline
	\includegraphics{anexos/ARCHI/application/event.png}         &
	\textbf{Application Event:} Evento que ocurre y afecta la continuidad de un proceso dentro del contexto de aplicaciones. \\
	\hline
	\includegraphics{anexos/ARCHI/application/object.png}        &
	\textbf{Data Object:} Objeto pasivo que es manipulado por comportamientos dentro del entorno de aplicaciones.                       \\
	\hline
\end{longtable}

\section{Elementos tecnológicos}

\begin{longtable}{|c|p{8cm}|}
	\caption{Elementos tecnológicos en ArchiMate} \label{tab:elementos-tecnologicos-archimate}                                                    \\
	\hline
	\textbf{Icono}                                              & \textbf{Descripción}                                                            \\
	\hline
	\endfirsthead

	\caption[]{Elementos tecnológicos en ArchiMate (continuación)}                                                                                \\
	\hline
	\textbf{Icono}                                              & \textbf{Descripción}                                                            \\
	\hline
	\endhead

	\hline
	\endfoot

	\endlastfoot
	\includegraphics{anexos/ARCHI/technology/facility.png}      &
	\textbf{Facility:} Espacio físico o infraestructura destinada al alojamiento de recursos tecnológicos. \\
	\hline
	\includegraphics{anexos/ARCHI/technology/equipment.png}     &
	\textbf{Equipment:} Recurso físico que puede ser utilizado en la ejecución de procesos organizacionales. \\
	\hline
	\includegraphics{anexos/ARCHI/technology/material.png}      &
	\textbf{Material:} Recursos físicos o insumos empleados o generados en el desarrollo de procesos empresariales. \\
	\hline
	\includegraphics{anexos/ARCHI/technology/node.png}          &
	\textbf{Node:} Recurso computacional encargado de alojar, procesar o interactuar con otros recursos y aplicaciones. \\
	\hline
	\includegraphics{anexos/ARCHI/technology/device.png}        &
	\textbf{Device:} Recurso físico capaz de ejecutar software de sistema o componentes de aplicación. \\
	\hline
	\includegraphics{anexos/ARCHI/technology/software.png}      &
	\textbf{System Software:} Software que proporciona la base necesaria para la ejecución de aplicaciones, como sistemas operativos o middleware. \\
	\hline
	\includegraphics{anexos/ARCHI/technology/collaboration.png} &
	\textbf{Technology Collaboration:} Interacción entre dos o más nodos o dispositivos que trabajan de manera conjunta. \\
	\hline
	\includegraphics{anexos/ARCHI/technology/interface.png}     &
	\textbf{Technology Interface:} Punto de acceso a través del cual un nodo o dispositivo ofrece sus funcionalidades. \\
	\hline
	\includegraphics{anexos/ARCHI/technology/process.png}       &
	\textbf{Technology Process:} Conjunto de comportamientos tecnológicos orientados al logro de un objetivo específico. \\
	\hline
	\includegraphics{anexos/ARCHI/technology/function.png}      &
	\textbf{Technology Function:} Agrupación de comportamientos tecnológicos relacionados entre sí. \\
	\hline
	\includegraphics{anexos/ARCHI/technology/interaction.png}   &
	\textbf{Technology Interaction:} Describe la interacción entre dos o más nodos o dispositivos. \\
	\hline
	\includegraphics{anexos/ARCHI/technology/service.png}       &
	\textbf{Technology Service:} Servicio proporcionado por la capa tecnológica a su entorno. \\
	\hline
	\includegraphics{anexos/ARCHI/technology/event.png}         &
	\textbf{Technology Event:} Evento que ocurre y afecta la continuidad de un proceso tecnológico. \\
	\hline
	\includegraphics{anexos/ARCHI/technology/artifact.png}      &
	\textbf{Artifact:} Objeto físico de información utilizado o generado en procesos tecnológicos. \\
	\hline
	\includegraphics{anexos/ARCHI/technology/network.png}       &
	\textbf{Communication Network:} Infraestructura que posibilita la comunicación entre nodos y dispositivos. \\
	\hline
	\includegraphics{anexos/ARCHI/technology/path.png}          &
	\textbf{Path:} Conexión que define el flujo de datos entre nodos o dispositivos. \\
	\hline
	\includegraphics{anexos/ARCHI/technology/distribution.png}  &
	\textbf{Distribution Network:} Red que facilita la distribución de productos o recursos a través de una infraestructura.                   \\
	\hline
\end{longtable}


\noindent
Si bien el lenguaje ArchiMate contempla capas adicionales como estrategia y relaciones, en el presente trabajo no se consideraron dichas capas, dado que el alcance del modelo se centra en la representación de los niveles de negocio, aplicación y tecnología.